\documentclass[11pt]{article}
\usepackage[margin=1in]{geometry}
\usepackage{amsmath,amssymb}
\usepackage[hidelinks]{hyperref}

\newcommand{\AP}{\mathrm{AP}}
\newcommand{\N}{\mathcal N}

\title{Literature-implied status for the nuclear-adjoint remark in arXiv:2204.12583}
\author{}
\date{2026-06-16}

\begin{document}
\maketitle

\paragraph{Status.}
This is a \emph{literature-implied answer}, not a new proof.  The source
remark is in Aksoy--Thiong, \emph{Schauder's Theorem and s-Numbers},
arXiv:2204.12583, Section ``Hutton's Theorem Revisited'' after the proof of
Hutton's theorem.  It says that, because nuclear operators are compact, it is
natural to ask how nuclearity of an operator \(T\) and of its adjoint \(T^*\)
are related.  The same remark recalls the positive theorem that if \(X^*\) has
the approximation property and \(T^*\) is nuclear, then \(T\) is nuclear and
\(\|T\|_{\mathcal N}=\|T^*\|_{\mathcal N}\).

\paragraph{Supporting identification.}
Reinov's paper \emph{On linear operators with \(s\)-nuclear adjoints:
\(0<s\le 1\)}, arXiv:1311.2270, gives the broader relationship.  Lemma 2
there states that for \(T\in L(X,Y)\),
\[
  \pi_Y T\in N_s(X,Y^{**}) \quad\Longleftrightarrow\quad
  T^*\in N_s(Y^*,X^*),
\]
and that \(T\in N_s(X,Y)\) implies \(T^*\in N_s(Y^*,X^*)\).  Theorem 1 states
that if \(s\in(0,1]\) and either \(X^*\) or \(Y^{***}\) has the approximation
property of order \(s\), then \(T^*\) being \(s\)-nuclear implies \(T\) is
nuclear.  For \(s=1\), this includes the usual nuclear-adjoint converse under
the classical approximation property assumptions.

\paragraph{Negative side.}
The converse fails in general without those approximation-property hypotheses.
Reinov's Theorem 2 states that for each \(r\in(2/3,1]\) there are a Banach
space \(Z_0\) and an operator \(T:Z_0^{**}\to Z_0\) such that \(Z_0^{**}\) has
a Schauder basis, \(Z_0^{***}\) has \(\AP_s\) for every \(s<r\),
\(\pi_{Z_0}T\in N_r(Z_0^{**},Z_0^{**})\), but
\[
  T\notin N_1(Z_0^{**},Z_0).
\]
Taking \(r=1\), Lemma 2 identifies \(\pi_{Z_0}T\in N_1\) with nuclearity of
the adjoint \(T^*\).  Thus there is a non-nuclear operator whose adjoint is
nuclear.  This supplies the negative general answer to the natural converse
question.

\paragraph{Scope.}
This answers only the nuclearity-versus-adjoint-nuclearity issue raised in the
remark.  It does not address the source paper's other themes: comparison of
general \(s\)-numbers, Hutton-type equalities, or Q-compactness under
approximation schemes.  Reinov's paper predates arXiv:2204.12583 and does not
explicitly cite that source, so this packet belongs in
\texttt{literature\_implied\_answers}, not \texttt{literature\_already\_answered}.

\begin{thebibliography}{9}
\bibitem{AT}
A. G. Aksoy and D. A. Thiong,
\emph{Schauder's Theorem and s-Numbers},
arXiv:2204.12583.

\bibitem{Reinov}
O. I. Reinov,
\emph{On linear operators with \(s\)-nuclear adjoints: \(0<s\le 1\)},
arXiv:1311.2270.
\end{thebibliography}

\end{document}

